Dependent Object Types (DOT) is the core calculus of Scala's type system. Its type safety is hard to prove because of bad bounds: under a binder of type $\dtyp{A}{\top}{\bot}$ every type is a subtype of every other, and the usual proof by preservation and progress needs to invert subtyping. This report explains a proof of DOT type safety that never inverts a DOT derivation. DOT typing derivations are translated into \FCdot, a calculus in which every use of subtyping is an explicit proof term that erases to nothing. Safety of \FCdot is store typing plus one theorem: over a typed store, closed evidence normalizes to a head form, by structural induction. Safety of DOT is then a corollary, for a store machine and for the reduction semantics of the WadlerFest presentation, with no restriction on recursive types.

The report is written for a reader who knows System F and the textbook proof of type safety but not DOT. It introduces DOT, shows where the textbook proof breaks, develops the idea on a fragment with functions and type members only, and then builds up the full calculus, its metatheory, and the translation. It also compares the proof with the direct proof of Rapoport et al., records the designs that did not work, and gives a tour of the \Lean mechanization, which has no library dependencies and depends on propositional extensionality and quotient soundness only.
